\section{Arhitectura sistemului}

\subsection{Componente logice}
\noindent
Sistemul are ca si componente principale:
{
\setlength{\parskip}{0pt}
\begin{enumerate}[noitemsep]
      \item \textbf{componenta de filtrare la nivel kernel}, responzabil de
            interceptarea pachetelor, extragerea informatiilor relevante despre acestea
            si transmiterea acestora, si blocarea pachetelor mentionate in fisierul de
            configuratii
      \item \textbf{componenta de intermediere}, responsabil de intermedierea
            informatiilor intre modulul kernel si serverul. Acesta de asemenea este
            responsabil si de citirea fisierelor de configuratii, sincronizarea acestora
            cu informatiile primite de la server si trimiterea noilor dispozitive ce
            trebuiesc blocate catre modul.
      \item \textbf{componenta centrala}, responsabil de stocarea tuturor
            informatiilor primite de la dispozitivele monitorizate si intermedierea
            legaturii intre totalitatea dispozitivelor si aplicatia adminului.
      \item \textbf{componenta de administrare}, ce ofera posibilitatea de
            vizualizare a conexiunilor ce au avut loc pe dispozitivele monitorizare si
            modificarea configuratiilor acestora intr-un mod individual sau global.
\end{enumerate}
}

\subsubsection{Componenta de filtrare la nivel kernel}
Aceasta componenta reprezinta un modul kernel scris in limbajul de programare C
ce ruleaza in spatiul adminului unitatii.  Principala caracteristica a acestuia
consta in definirea si inregistrarea unui \textit{`hook'} utilizand frameworkul
\textit{`Netfilter'}, care permite interceptarea, analiza si blocarea pachetelor
in retea, in anumite puncte specifice ale procesului de procesare.

Cu ajutorul acestui framework este registrata pentru reapelare functia
\textit{`packet\_hook'} la intrarea in procesarea locala
(\textit{`NF\_INET\_LOCAL\_IN'}) a pachetelor.

\subsubsection{Componenta de intermediere}
Aceasta componenta reprezinta un agent local, \textit{`daemon'}, ce ruleaza un
spatiul utilizatorului. Limbajul ales pentru crearea acestei componente a fost
C++, pentru a utiliza componente cu un grad de abstractizare mai mare, pastrand
totodata puterea si viteza de procesare a limbajului C. Agentul este responsabil
de trimiterea informatiilor ce au fost extrase din pachete in modulul kernel
spre unitatii centrale, si de asemenea de sincronizarea setarilor si regulilor
noi selectate in aplicatie cu regulile si modul de functionare a componentei de
filtrare.

Programul comunica cu modulul kernel este facuta prin intermediul unui
\textit{Miscellaneous character device}, ce ofera posibilitatea folosirii unui
\textit{character device} intr-un mod simplificat. Modulul kernel trimite
informatii referitoare la pachetele primite intr-un \textit{character device},
informatii ce sunt periodic citite de agentul local si trimise in grupari catre
componenta centrala.

Intr-un mod similar, programul trimite periodic mesaje catre server pentru a
verifica daca configurarea aplicatiei trebuie modificata. In cazul in care
serverul il informa ca modificari au avut loc, acesta va scrie in fisierul de
tip \textit{character device} pentru a modifica modul de functionare a
componentei de filtrare.

% revormulez partea despre 3 fisiere 
De asemenea, componenta este responsabila de stocarea locala a configuratiei
primite de la server, cat si de citirea acesteia la inceputul rularii
programului. Acesta stocare se face in 3 fisiere, responsabile pentru stocarea
unei indexari unice a dispozitivului, salvarea modului de rulare a aplicatiei, a
locatiei fisierului ce contine lista de conexiuni blocate si informatii despre
conexiunea cu serverul (adresa ip si portul ce trebuie utilizat), si fisierul
responsabil de retinerea dispozitivelor ce trebuiesc blocate.

\subsubsection{Componenta centrala}
Componenta centrala reprezinta serverul aplicatiei responsabila de gestionarea
legaturii dintre dispozitivele utilizate si aplicatia desktop. Aceasta este
responsabila de centralizarea informatiilor transmise de modulele kernel ce
ruleaza pe dispozitivele in legatura directa cu acesta, si reformatarea acesteia
pentru a putea fi vizualizata si prelucrata de aplicatia utilizatorului admin.
Totodata, prin componenta centrala, noile reguli setate din aplicatia desktop sunt
transmise catre dispozitive.

\subsubsection{Componenta de administrare}
Componenta de administrare este reprezentata de aplicatia desktop ce ofera
flexibilitate utilizatorilor pentru control asupra dispozitivelor asocitate
serverului. Prin aceasta un utilizator poate modifica atat modul de rulare, cat
si reguli de blocare individuala a conexiunilor externe ce au fost, sau pot fi
in legatura cu dispozitivele. De asemenea, ofera posibilitatea de vizualizare de
tip graf a conexiunilor care au fost facute cu dispozitivele suport.